% Source: https://artofproblemsolving.com/community/c3h1586497_remaining_surface_area
% Stable identity: aops-topic-1586497-q01
The remaining solid is bounded by the curved surface of the cylinder and by the two hollow cone surfaces; the two circular ends are gone, because they were the bases of the removed cones.

\begin{asy}
usepackage("amsmath");
size(7cm);
defaultpen(fontsize(10pt));
real L=8,r=3,k=0.35;
path el(real x){return shift(x,0)*xscale(k*r)*yscale(r)*unitcircle;}
fill((0,r)--(L/2,0)--(0,-r)--cycle,gray(0.85));fill((L,r)--(L/2,0)--(L,-r)--cycle,gray(0.85));
draw((0,r)--(L,r),linewidth(0.9));draw((0,-r)--(L,-r),linewidth(0.9));
draw(el(0),linewidth(0.9));draw(el(L),linewidth(0.9));
draw((0,r)--(L/2,0)--(0,-r),dashed);draw((L,r)--(L/2,0)--(L,-r),dashed);
draw((0,0)--(L/2,0),dotted);draw((L/2,0)--(L,0),dotted);
dot((L/2,0));
label("$5$",(L/4,r/2),NE);label("$4$",(L/4,0),S);
draw((0,-r-0.7)--(L,-r-0.7),Arrows(TeXHead));label("$8$",(L/2,-r-0.7),S);
draw((-1.6,-r)--(-1.6,r),Arrows(TeXHead));label("$6$",(-1.6,0),W);
label("\mbox{}",(min(currentpicture,user=true)+max(currentpicture,user=true))/2,black); // keeps the text after the figure black when LaTeX typesets the labels inline
\end{asy}

\textbf{Dimensions.} The radius is $3$ cm and each cone has height $\dfrac{8}{2}=4$ cm, so its slant height is $\sqrt{3^2+4^2}=5$ cm.

\textbf{Area.}
\[
\begin{aligned}
\text{curved surface of the cylinder}&=2\pi rh=2\pi\cdot3\cdot8=48\pi,\\
\text{curved surface of each cone}&=\pi r\ell=\pi\cdot3\cdot5=15\pi .
\end{aligned}
\]
The total is $48\pi+2\cdot15\pi=78\pi\approx78\cdot3.142=\boxed{245.08\text{ cm}^2}$ (to two decimal places).

\emph{Why subtracting fails.} Removing a cone does not remove its lateral surface from the total; it exposes it as a new hollow surface of the solid. Only the cylinder's two ends disappear. So ``surface area of the cylinder minus surface area of the cones'' is not the surface area of what remains.
